% Section: system design of NetworkLogger (app + pipeline + APIs).
% Compile inside a report/article class; needs no extra packages.

\section{System Design}

\subsection{Android Application}
NetworkLogger is a Kotlin Android application that samples the mobile-network
environment every 5 seconds. A foreground service (\texttt{LoggingService})
runs the sampling loop under a partial wake lock so that logging survives the
screen turning off. Each sample records the serving cell (PCI, TAC/eNB, cell
ID), signal quality (RSRP, RSRQ, SINR), the LTE/NR band, the GPS position, and
passive uplink/downlink throughput. Samples are buffered on the device and
uploaded once per minute as a batch (\texttt{CloudSync}), with an offline queue
so that data collected without connectivity is uploaded later. A periodic
15-minute worker (\texttt{NetworkWorker}) logs device health (battery, storage)
and acts as a watchdog that restarts logging if the service was killed; killed
sessions resume under the same session identifier, and every session keeps an
event log of why logging stopped (\texttt{SessionDiagnostics}). The app ships
in two build flavors: an \emph{internal} flavor with a developer mode that
keeps CSV files on the phone, and a \emph{public} flavor for volunteer testers
that always streams to the cloud.

\subsection{Data Pipeline}
Data flows from the phones to dashboards as follows:
\[
\text{Phone} \;\rightarrow\; \text{Firestore} \;\rightarrow\;
\text{listener (Python)} \;\rightarrow\; \text{InfluxDB} \;\rightarrow\;
\text{Grafana}
\]
Each phone writes its batches into Cloud Firestore under a per-user,
per-session document hierarchy
(\texttt{users/\{uid\}/sessions/\{session\}/batches/\{seq\}}). A Python
listener (\texttt{firestore\_listener.py}) subscribes to Firestore in real
time and writes every 5-second row into InfluxDB, a time-series database with
a 180-day retention policy. Grafana reads from InfluxDB and renders two live
dashboards: a time-series ``Live'' dashboard (signal quality, throughput,
gaps) and a geographic ``Map'' dashboard. The pipeline runs both on a
development machine and on a lab server, where a cron watchdog keeps the
listener alive and the dashboards are shared view-only on the institute
network. A separate offline toolchain (Python, pandas/matplotlib) exports the
Firestore data and produces analysis charts, e.g.\ on logging robustness and
per-device comparisons.

\subsection{APIs Used}
\begin{itemize}
  \item \textbf{TelephonyManager / CellInfo APIs}
        (\texttt{android.telephony}): serving-cell identity and measurements
        --- \texttt{CellInfoLte}/\texttt{CellInfoNr},
        \texttt{CellIdentityLte}/\texttt{CellIdentityNr} (PCI, cell ID, EARFCN
        / NR-ARFCN, from which the band is derived), and
        \texttt{SignalStrength} / \texttt{CellSignalStrengthNr} via a
        \texttt{PhoneStateListener} (RSRP, RSRQ, SINR, network type,
        operator).
  \item \textbf{Fused Location Provider}
        (\texttt{com.google.android.gms.location}): high-accuracy GPS fixes
        for geotagging every sample.
  \item \textbf{TrafficStats} (\texttt{android.net.TrafficStats}): passive
        throughput, computed from the difference in transmitted/received byte
        counters between consecutive samples (excluding the app's own
        traffic).
  \item \textbf{PowerManager} (\texttt{android.os.PowerManager}): a partial
        wake lock that keeps the CPU running during screen-off logging.
  \item \textbf{WorkManager} (\texttt{androidx.work}): the periodic
        device-health logger and watchdog that revives killed sessions.
  \item \textbf{BatteryManager / StatFs} (\texttt{android.os}): battery level
        and free-storage readings for the device-health records.
  \item \textbf{Firebase Authentication}
        (\texttt{com.google.firebase.auth}): anonymous sign-in that gives each
        device a stable user ID used in the Firestore layout and security
        rules.
  \item \textbf{Cloud Firestore} (\texttt{com.google.firebase.firestore}):
        batched uploads of samples, session metadata and status, with
        built-in offline persistence.
  \item \textbf{Firestore Admin SDK (Python)}: server-side real-time listener
        that moves new batches into the time-series store.
  \item \textbf{InfluxDB client API (Python)}: writes the samples as
        time-series points and serves Flux queries to Grafana.
  \item \textbf{Grafana HTTP API / dashboards}: dashboard JSON is generated by
        script and provisioned on both the development machine and the lab
        server.
\end{itemize}
